\chapter{FPGA 엔진의 내부: \texorpdfstring{\code{bombe\_engine}}{bombe\_engine}}\label{ch:fpga-engine}

기준 RTL은 칩에 올릴 수 없다(\ref{sec:sim-vs-synth}절). 실제 FPGA에 올리는 것은 같은 동작을 칩에 맞는
구조로 다시 쓴 \emph{FPGA 엔진}(\code{rtl/fpga/src/bombe\_engine.sv})이다. 엔진은 이 프로젝트의 세 번째
봄베이기도 하다. 기어로 도는 기계식, 릴레이식과 나란히 놓여 같은 메뉴에서 같은 정지를 내는, 움직이는
부품이 없는 전자식이다.

엔진 하나는 로터 순서 하나의 모든 드럼 위치를 훑는다. 기준 RTL의 \code{bombe\_core}와 \code{bombe\_sweep}을
합친 일이다. F부에서 개념으로만 소개한 것들(바뀌는 것만 다시 계산하기, 파이프라인, 분산 RAM, 케이블 단위
전파, FIFO)이 이 한 파일에 모두 들어 있다. 이 장은 그 코드를 신호가 흐르는 순서대로 따라간다.
그림~\ref{fig:g-engine-blocks}이 전체 구성이다.

\begin{figure}[ht]
\centering
\begin{tikzpicture}[x=1cm,y=1cm]
  \node[blkg,minimum width=22mm,minimum height=14mm] (host) at (-0.2,-0.1) {호스트\\{\scriptsize PC 또는 보드의 ARM}};
  \node[blkn,minimum width=22mm,minimum height=14mm] (br) at (2.8,-0.1) {브리지\\{\scriptsize UART / AXI4-Lite}};
  \draw[rounded corners=2pt,draw=cgray,fill=clight] (5.9,-3.1) rectangle (14.5,2.0);
  \node[anchor=north west,font=\small\bfseries] at (6.0,1.95) {\code{bombe\_array}: 엔진 $N$대};
  \draw[rounded corners=2pt,draw=cblue,fill=white] (6.2,-2.7) rectangle (14.2,1.2);
  \node[anchor=north west,font=\small\bfseries] at (6.3,1.15) {\code{bombe\_engine}};
  \node[blko,minimum width=21mm] (cfg) at (7.5,-0.1) {설정 레지스터\\{\scriptsize 메뉴, 로터 순서}};
  \node[blkb,minimum width=21mm] (tbl) at (10.2,-0.1) {치환표 계산\\{\scriptsize 파이프라인}};
  \node[reg,minimum width=20mm] (ram) at (12.8,-0.1) {표 메모리\\{\scriptsize 분산 RAM}};
  \node[blkb,minimum width=21mm] (spr) at (10.2,-1.8) {전파\\{\scriptsize 676비트 상태}};
  \node[blkr,minimum width=21mm] (fifo) at (7.5,-1.8) {정지 FIFO};
  \draw[->] (cfg) -- (tbl); \draw[->] (tbl) -- (ram);
  \draw[->] (ram.south) |- (spr.east);
  \draw[->] (spr) -- (fifo);
  \draw[<->] (host) -- (br);
  \draw[<->] (br) -- node[above,font=\scriptsize,text=cgray] {레지스터 버스} (6.2,-0.1);
\end{tikzpicture}
\caption{FPGA 엔진의 블록 수준 구성. 이 장은 파란 테두리 안을, 다음 장은 그 바깥을 다룬다.}
\label{fig:g-engine-blocks}
\end{figure}

\section{바뀌는 것만 다시 계산한다}\label{sec:e-recompute}

\Why 계산의 입력이 바뀌지 않았으면 결과도 바뀌지 않는다. 그러니 입력이 \emph{얼마나 자주} 바뀌는지에 따라
계산을 층으로 나누고, 드물게 바뀌는 층의 결과는 저장해 두었다가 다시 쓰면 된다. 회로에서 이 생각의 값어치가
특히 큰 까닭은 \ref{sec:tradeoff}절에서 보았다. 자주 하지 않는 계산은 작은 회로로 여러 클럭에 걸쳐 천천히
해도 전체 시간에 거의 영향을 주지 않는다.

\General 자동차의 주행 거리계를 보자. 맨 오른쪽 숫자 바퀴는 계속 돌지만 그 왼쪽 바퀴는 열 번에 한 번,
그 왼쪽은 백 번에 한 번 움직인다. 숫자판 전체를 매번 새로 그리는 사람은 없다. 회로의 예로는 디지털 시계의
표시부가 있다. ``몇 월 며칠 무슨 요일''을 계산하는 회로는 복잡하지만 하루에 한 번만 돌리면 되므로, 덧셈기
하나를 돌려 쓰는 느린 회로로 만들고 결과를 레지스터에 넣어 둔다. 초를 표시하는 쪽만 매번 움직인다.

\Ours 스크램블러 하나의 치환은 식으로 쓰면 다음과 같다. 전류는 오른쪽 로터를 지나($R_k$), 가운데와 왼쪽
로터와 반사판을 돌아 나와($I$), 다시 오른쪽 로터를 거꾸로 지난다($R_k^{-1}$).
\[
  S_k \;=\; R_k^{-1} \circ I \circ R_k, \qquad
  I \;=\; M^{-1} \circ L^{-1} \circ U \circ L \circ M
\]
$k$는 오른쪽 로터의 위치이고, $I$가 \ref{sec:tradeoff}절에서 말한 ``안쪽''이다(4로터 기종에서는 $L$과 $U$
사이에 그리스 로터가 낀다). 오른쪽 로터가 한 칸 갈 때 바뀌는 것은 $k$뿐이고 $I$는 그대로다. 그리고 $k$가
가질 수 있는 값은 26가지뿐이므로, $I$가 정해지면 $S_0$부터 $S_{25}$까지를 한꺼번에 표로 만들어 둘 수 있다.

엔진은 계산을 세 층으로 나눈다(그림~\ref{fig:g-e-layers}).

\begin{figure}[ht]
\centering
\begin{tikzpicture}[x=1cm,y=1cm]
  \draw[rounded corners=2pt,draw=cgray,fill=clight] (0,0) rectangle (15.4,5.0);
  \node[anchor=north west,font=\small\bfseries] at (0.15,4.95) {구동에 한 번: 고른 로터의 배선표를 레지스터로 옮긴다 (\code{LOAD}, 234클럭)};
  \draw[rounded corners=2pt,draw=corange,fill=corangebg] (0.4,0.3) rectangle (15.0,4.1);
  \node[anchor=north west,font=\small\bfseries] at (0.55,4.05) {가운데 로터가 바뀔 때 (676번): 안쪽 치환 $I$와 표 $S_0\ldots S_{25}$를 만든다};
  \node[anchor=north west,font=\footnotesize,text=cgray] at (0.55,3.5) {\code{INNER} 약 42클럭 + \code{SROWS} 약 33클럭};
  \draw[rounded corners=2pt,draw=cblue,fill=cbluebg] (0.8,0.6) rectangle (14.6,2.8);
  \node[anchor=north west,font=\small\bfseries] at (0.95,2.75) {위치마다 (17,576번): 표에서 읽고 전파한다};
  \node[anchor=north west,font=\footnotesize,text=cgray] at (0.95,2.2) {\code{LATCH} 1클럭 + \code{SPREAD} 여러 클럭 + \code{ADVANCE} 1클럭};
  \node[anchor=north west,font=\footnotesize,text=cgray] at (0.95,1.6) {정지이면 \code{PUSH} 1클럭을 더해 FIFO에 넣는다};
\end{tikzpicture}
\caption{엔진의 세 층. 횟수는 3로터 기종의 로터 순서 하나에 대한 것이다. 바깥 층일수록 드물게 돌고, 드물게
도는 층일수록 여러 클럭을 써도 된다.}
\label{fig:g-e-layers}
\end{figure}

가운데 층은 한 번에 약 75클럭이 들지만 26위치에 한 번이므로 위치 하나로 나누면 3클럭이 못 된다. 그 대가로
기준 RTL에서 메뉴의 선마다 하나씩(최대 16벌) 있던 스크램블러 회로가 계산 회로 한 벌과 메모리로 바뀐다.
실제 기계와도 같은 모양이다. 봄베의 느린 드럼 둘은 빠른 드럼이 한 바퀴 도는 동안 가만히 있다.

이 세 층을 차례로 돌리는 것이 엔진의 상태 기계다(그림~\ref{fig:g-e-fsm}). 상태가 여덟이고, 상태마다 이 장의
한 절이 대응한다.

\begin{figure}[ht]
\centering
\begin{tikzpicture}[x=1cm,y=1cm,
   st/.style={draw,rounded corners=3pt,minimum width=19mm,minimum height=8mm,fill=cbluebg,draw=cblue,font=\small\ttfamily}]
  \node[st,fill=clight,draw=cgray] (idle) at (0,0) {IDLE};
  \node[st,fill=clight,draw=cgray] (load) at (3.0,0) {LOAD};
  \node[st,fill=corangebg,draw=corange] (inner) at (6.0,0) {INNER};
  \node[st,fill=corangebg,draw=corange] (srows) at (9.0,0) {SROWS};
  \node[st] (latch) at (12.0,0) {LATCH};
  \node[st] (spread) at (12.0,-2.2) {SPREAD};
  \node[st,fill=credbg,draw=cred] (push) at (8.0,-2.2) {PUSH};
  \node[st] (adv) at (4.0,-2.2) {ADVANCE};
  \draw[->] (idle) -- node[above,note] {START} (load);
  \draw[->] (load) -- (inner); \draw[->] (inner) -- (srows); \draw[->] (srows) -- (latch);
  \draw[->] (latch) -- (spread);
  \draw[->] (spread) -- node[above,note] {정지} (push);
  \draw[->] (push) -- (adv);
  \draw[->] (spread.south) -- ++(0,-0.7) -| node[below,note,pos=0.25] {통과 (테스트 레지스터가 가득 참)} (adv.south);
  \draw[->] (adv.north) -- ++(0,0.95) -| node[above,note,pos=0.25] {가운데 로터가 바뀜} (inner.south);
  \draw[->] ([xshift=5mm]adv.north) -- ++(0,0.45) -| node[above,note,pos=0.3] {오른쪽 로터만 바뀜} ([xshift=-4mm]latch.south);
  \draw[->] (adv.west) -| node[below,note,pos=0.25] {마지막 위치} (idle.south);
\end{tikzpicture}
\caption{엔진의 상태 기계. 회색은 구동에 한 번, 주황은 가운데 로터가 바뀔 때, 파랑은 위치마다 지나는 상태다.}
\label{fig:g-e-fsm}
\end{figure}

\section{배선표 적재기}\label{sec:e-loader}

\begin{lstlisting}[language=SV,caption={\code{bombe\_engine.sv}: 배선표를 한 칸씩 옮기는 두 단짜리 적재기.},label={lst:e-loader}]
logic [129:0] tbl [9];     // 9 tables x 26 letters x 5 bits

always_ff @(posedge clk) begin
  ld_v1   <= (state == S_LOAD);
  ld_t1   <= ld_t;
  ld_c1   <= ld_c;
  ld_src1 <= (ld_t == 4'd8) ? 2'd2 : {1'b0, ld_t[0]};
  ld_idx1 <= ({7'd0, ld_id} * 11'd26 + {6'd0, ld_c}) * 11'd5;
  if (ld_v1) tbl[ld_t1][ld_c1*5 +: 5] <= ld_val;
end

always_comb begin
  case (ld_src1)
    2'd0:    ld_val = wirings_pkg::ROTOR_FWD[ld_idx1 +: 5];
    2'd1:    ld_val = wirings_pkg::ROTOR_BWD[ld_idx1 +: 5];
    default: ld_val = wirings_pkg::REFLECTOR[ld_idx1[9:0] +: 5];
  endcase
end
\end{lstlisting}

\Why 로터의 배선표는 상수이고 \code{wirings\_pkg}에 로터 열 개 치가 한 줄로 이어져 있다(\ref{ch:rtl-scrambler}장).
어느 로터를 쓸지는 호스트가 실행 중에 정한다. 그러면 표를 읽는 곳마다 ``열 개의 로터 가운데 하나, 그
안의 26칸 가운데 하나''를 고르는 260:1 선택 회로가 생긴다(\ref{sec:mux-cost}절). 표를 읽는 곳은
파이프라인의 여러 단과 26개의 레인에 흩어져 있으므로 이 큰 선택 회로가 수십 벌 복사된다. 그런데 로터의
번호는 구동 중에 바뀌지 않는다. 구동을 시작할 때 고른 로터의 표만 작은 레지스터로 \emph{옮겨 두면},
그 뒤의 모든 조회는 26칸 가운데 하나를 고르는 26:1 선택으로 줄어든다.

\General 요리사는 주문이 들어올 때마다 창고에 가지 않는다. 영업 전에 그날 쓸 재료만 조리대에 꺼내 둔다.
회로의 예로는 통신 속도를 스위치로 고르는 직렬 통신기가 있다. 속도마다 다른 분주값을 표에서 읽어야 하는데,
전원을 켤 때 한 번 읽어 레지스터에 넣어 두고 그 뒤로는 레지스터만 본다. 표를 고르는 회로는 한 벌이면 되고
빠를 필요도 없다.

\Ours \code{tbl}이 그 조리대다. 표가 아홉 개다. 오른쪽, 가운데, 왼쪽, 그리스 로터 각각의 정방향과 역방향
표 여덟 개와 반사판 표 하나이고, 표 하나는 26칸 $\times$ 5비트 = 130비트다. \code{LOAD} 상태에서 계수기 둘
(\code{ld\_t}는 표 번호, \code{ld\_c}는 칸 번호)이 $9 \times 26 = 234$칸을 한 클럭에 한 칸씩 지나간다.

적재기는 두 단이다. 첫 단은 상수 표 안에서의 비트 위치 $(26 \times \text{로터 번호} + \text{칸}) \times 5$를
계산해 \code{ld\_idx1}에 담는다. 둘째 단은 그 위치의 5비트를 읽어(\code{ld\_val}) \code{tbl}의 해당 칸에
써넣는다. 이것을 한 단으로 하면 곱셈과 덧셈으로 위치를 만들고, 1,300비트짜리 상수에서 5비트를 골라내고,
그것을 써넣을 칸을 해독하는 일이 한 클럭에 이어진다. 가장 긴 경로가 될 까닭이 없는 곳에서 타이밍을 잡아먹지
않도록 가운데에 레지스터를 하나 넣어 둘로 잘랐다(\ref{sec:pipeline}절의 방법 그대로다). 그림~\ref{fig:g-e-loader}의 파형처럼 칸마다 한 클럭씩
늦게 써질 뿐 한 클럭에 한 칸이라는 속도는 그대로다.

\begin{figure}[ht]
\centering
\begin{tikztimingtable}[timing/slope=0.1,timing/coldist=4pt,xscale=2.4,yscale=1.25,semithick]
  \code{clk}          & 7{.5L .5H} \\
  \code{ld\_c}        & .5D{0} 1D{1} 1D{2} 1D{3} 1D{4} 1D{5} 1D{6} .5D{7} \\
  \code{ld\_idx1}     & .5U 1D{0} 1D{1} 1D{2} 1D{3} 1D{4} 1D{5} .5D{6} \\
  \code{tbl}에 써진 칸 & 1.5U 1D{0} 1D{1} 1D{2} 1D{3} 1D{4} .5D{5} \\
\extracode
  \begin{pgfonlayer}{background}
    \vertlines[cgray!50,dashed]{0.5,1.5,...,6.5}
  \end{pgfonlayer}
\end{tikztimingtable}
\caption{적재기의 두 단. 계수기가 가리킨 칸은 한 클럭 뒤에 위치가 계산되고 두 클럭 뒤에 \code{tbl}에 들어 있다.}
\label{fig:g-e-loader}
\end{figure}

마지막 칸은 \code{LOAD} 상태가 끝난 뒤에야 써진다. 그래도 괜찮은 이유는 마지막에 옮기는 표가 반사판이고,
바로 다음 상태인 \code{INNER}의 파이프라인에서 첫 글자가 반사판에 닿는 것은 일곱 단 뒤이기 때문이다.
``늦게 도착하지만 쓰이기 전에는 도착한다''는 것을 따져 보고 기다리는 상태를 따로 두지 않았다.

코드의 \code{ld\_id}는 지금 옮기는 표가 어느 로터의 것인지를 설정 레지스터에서 고른 값이다. 3로터 기종에서는
그리스 로터의 번호에 쓰레기 값이 들어 있을 수 있으므로, 범위를 벗어난 번호는 0으로 바꿔 상수 표 바깥을 읽지
않게 한다. 그렇게 읽은 표는 어차피 쓰이지 않는다.

\section{안쪽 치환 파이프라인}\label{sec:e-inner}

안쪽 치환 $I$는 글자 하나가 가운데 로터, 왼쪽 로터, (그리스 로터,) 반사판을 지나 같은 길로 돌아오는
계산이다. \ref{sec:pipeline}절에서 이것이 열네 단 파이프라인이 된 사정을 보았다. 여기서는 코드를 본다.

\begin{lstlisting}[language=SV,caption={\code{bombe\_engine.sv}: 안쪽 치환 파이프라인(가운데 로터의 정방향과 끝부분만).},label={lst:e-inner}]
always_ff @(posedge clk) begin
  in_valid <= {in_valid[IN_ST-2:0], in_feed};
  in_xs[0] <= in_x;
  in_va[0] <= in_x;
  in_vb[0] <= in_x;
  for (int s = 1; s < IN_ST; s++) in_xs[s] <= in_xs[s-1];
  // middle forward
  in_va[1]  <= wirings_pkg::add26(in_va[0], m);
  in_vb[1]  <= wirings_pkg::add26(in_vb[0], m1_q);
  in_va[2]  <= wirings_pkg::sub26(pick(tbl[2], in_va[1]), m);
  in_vb[2]  <= wirings_pkg::sub26(pick(tbl[2], in_vb[1]), m1_q);
  // ... left forward, greek forward, reflector, greek backward, left backward ...
  // middle backward
  in_va[12] <= wirings_pkg::add26(in_va[11], m);
  in_va[13] <= wirings_pkg::sub26(pick(tbl[3], in_va[12]), m);
  if (in_valid[IN_ST-1]) begin
    inner0[in_xs[IN_ST-1]*5 +: 5] <= in_va[IN_ST-1];
    inner1[in_xs[IN_ST-1]*5 +: 5] <= in_vb[IN_ST-1];
  end
end
\end{lstlisting}

\paragraph{로터 한 단에 두 단.} 로터를 지나는 계산은 \ref{sec:through}절의 \code{through}와 같다. 글자에
로터의 회전량을 더하고, 배선표를 읽고, 회전량을 다시 뺀다. 코드에서 \code{in\_va[1]}이 더하기이고
\code{in\_va[2]}가 읽고 빼기다. \code{pick(tbl[2], c)}는 130비트 표에서 \code{c}번째 5비트를 꺼내는 함수다.

\Why 왜 이 둘을 한 단에 넣지 않았는가. 한 단의 길이는 신호가 지나는 논리의 양이 정한다. 26을 법으로 하는
덧셈(\code{add26})은 그냥 덧셈이 아니다. 더하고, 26 이상인지 비교하고, 넘었으면 26을 뺀다. 표 읽기는 5비트
폭의 26:1 선택이다. 뺄셈(\code{sub26})도 덧셈과 같은 크기다. 이 셋을 이으면 덧셈기 둘과 선택 회로 하나가
한 줄이 되고, 여기에 100MHz 클럭의 10나노초는 모자랐다. 덧셈과 ``읽고 빼기'' 사이에 레지스터를 두면 각
단은 그 절반쯤이 된다. 단이 일곱에서 열넷으로 늘어 결과가 일곱 클럭 늦게 나오지만, 이 계산은 26위치에
한 번 하므로 위치 하나에 0.3클럭도 안 되는 값이다.

\General 공장의 조립 줄에서 한 작업자가 ``나사를 고르고, 끼우고, 조이는'' 일을 다 하면 줄 전체가 그 사람의
속도로 움직인다. 고르는 사람과 조이는 사람을 나누면 줄은 빨라지고, 제품 하나가 줄을 다 지나는 시간만 조금
늘어난다. 회로의 예로는 곱셈기가 있다. 큰 곱셈을 한 클럭에 하지 않고 부분곱을 만드는 단과 그것을 더하는
단으로 나누는 것이 보통이다.

\paragraph{꼬리표를 함께 흘린다.} 파이프라인에는 글자 여러 개가 동시에 들어 있다. 맨 끝에서 결과가 나왔을
때 그것이 어느 글자의 결과인지 알아야 표의 맞는 칸에 써넣을 수 있다. 그래서 원래 글자 \code{in\_x}를
\code{in\_xs}라는 레지스터 줄에 실어 결과와 나란히 흘려보낸다. \code{for} 문 한 줄이 그것이다(반복문이므로
레지스터 열세 개가 한 줄로 이어진 회로가 된다). \code{in\_valid}도 같다. ``이 단에 진짜 글자가 들어 있다''는
한 비트를 단마다 하나씩 옮긴다. 맨 끝 단의 비트가 1일 때에만 표에 쓴다(그림~\ref{fig:g-e-inner}).

\begin{figure}[ht]
\centering
\begin{tikzpicture}[x=0.86cm,y=1cm]
  \node[anchor=east,font=\footnotesize] at (-0.2,1.6) {유효 \code{in\_valid}};
  \node[anchor=east,font=\footnotesize] at (-0.2,0.8) {꼬리표 \code{in\_xs}};
  \node[anchor=east,font=\footnotesize] at (-0.2,0.0) {값 \code{in\_va}};
  \foreach \i/\pv/\px/\pa in {0/1/D/D,1/1/C/{+m},2/1/B/M,3/1/A/{+l},4/0/{}/{}} {
    \draw[fill=cyellowbg] (\i*1.6,1.3) rectangle ++(1.1,0.6); \node[font=\ttfamily\footnotesize] at (\i*1.6+0.55,1.6) {\pv};
    \draw[fill=cyellowbg] (\i*1.6,0.5) rectangle ++(1.1,0.6); \node[font=\ttfamily\footnotesize] at (\i*1.6+0.55,0.8) {\px};
    \draw[fill=cbluebg,draw=cblue] (\i*1.6,-0.3) rectangle ++(1.1,0.6); \node[font=\footnotesize] at (\i*1.6+0.55,0.0) {\pa};
    \node[font=\tiny,text=cgray] at (\i*1.6+0.55,-0.6) {\i 단};
  }
  \foreach \i in {0,1,2,3} { \draw[->] (\i*1.6+1.1,1.6) -- ++(0.5,0); \draw[->] (\i*1.6+1.1,0.8) -- ++(0.5,0); \draw[->] (\i*1.6+1.1,0.0) -- ++(0.5,0); }
  \node at (8.5,0.8) {$\cdots$};
  \draw[fill=cyellowbg] (9.2,1.3) rectangle ++(1.1,0.6); \node[font=\ttfamily\footnotesize] at (9.75,1.6) {1};
  \draw[fill=cyellowbg] (9.2,0.5) rectangle ++(1.1,0.6); \node[font=\ttfamily\footnotesize] at (9.75,0.8) {Q};
  \draw[fill=cbluebg,draw=cblue] (9.2,-0.3) rectangle ++(1.1,0.6); \node[font=\footnotesize] at (9.75,0.0) {$I(\mathtt{Q})$};
  \node[font=\tiny,text=cgray] at (9.75,-0.6) {13단};
  \node[reg,minimum width=24mm,minimum height=14mm] (tb) at (12.9,0.4) {\code{inner0}\\{\scriptsize 26칸}};
  \draw[->] (10.3,0.0) -- node[below,note] {값} (tb.west|-0,0.0);
  \draw[->] (10.3,0.8) -- node[above,note] {칸 번호} (tb.west|-0,0.8);
  \draw[->] (10.3,1.6) -| node[above,note,pos=0.3] {쓰기 허가} (tb.north);
\end{tikzpicture}
\caption{파이프라인의 세 줄. 값과 함께 ``누구의 값인가''(꼬리표)와 ``진짜 값인가''(유효 비트)가 같은 속도로
흐른다. 글자 A, B, C, D가 차례로 들어간 뒤의 한 순간이다.}
\label{fig:g-e-inner}
\end{figure}

\General 택배 상자에 붙은 송장이 꼬리표다. 컨베이어 위의 상자는 내용물만으로는 어디로 갈지 알 수 없다.
처리 장치의 파이프라인도 같은 일을 한다. 계산 결과와 함께 ``이 결과를 몇 번 레지스터에 쓸 것인가''를 단마다
옮긴다.

\paragraph{두 줄을 나란히.} 코드에는 \code{in\_va}와 \code{in\_vb}가 짝으로 나온다. \code{in\_va}는 가운데
로터의 위치를 \code{m}으로, \code{in\_vb}는 \code{m}보다 한 칸 넘어간 \code{m1\_q}로 계산한다. 결과는 각각
\code{inner0}($I$)과 \code{inner1}($I'$)에 들어간다. $I'$은 크립의 중간에서 가운데 로터가 넘어갔다고 가정하는
경우(\ref{ch:rust-bombe}장의 넘어감 가설)에, 넘어간 뒤쪽의 선들이 쓰는 표다. 같은 회로를 한 벌 더 둔 것이고
조회하는 표는 같다. \code{m1\_q}는 $m + 1$을 미리 레지스터에 담아 둔 것이다. 파이프라인의 여러 단이 이 값을
쓰는데, 쓸 때마다 더하면 단마다 덧셈이 하나씩 길어진다.

\paragraph{그리스 로터가 없을 때.} 3로터 기종에서는 그리스 로터의 네 단이 값을 그대로 넘긴다
(\code{greek\_en ? ... : in\_va[4]}). 단을 건너뛰지 않고 빈 단으로 남겨 둔 것은 단 수가 기종에 따라
달라지면 제어가 복잡해지기 때문이다. 단 수가 언제나 열넷이면 ``26글자를 넣고 열네 클럭을 기다린다''는
제어 하나로 끝난다. 제어는 이렇다. \code{INNER} 상태에서 \code{in\_x}를 0부터 25까지 올리며 한 클럭에 한
글자씩 넣고, 다 넣은 뒤 \code{wait\_cnt}로 단 수만큼 기다려 마지막 글자가 빠져나오게 한다. 합쳐서 약 42클럭이다.

\section{26레인 행 만들기}\label{sec:e-rows}

안쪽 치환이 준비되면 오른쪽 로터의 26위치 각각에 대한 스크램블러 치환 $S_k = R_k^{-1} \circ I \circ R_k$를
만든다. $S_k$ 하나는 26글자의 치환이므로 130비트 한 \emph{행}이고, 행이 26개다(넘어감 가설을 쓰면 $I'$으로
만든 26행이 더해져 52개다).

\begin{lstlisting}[language=SV,caption={\code{bombe\_engine.sv}: 한 클럭에 한 행을 만드는 26레인, 5단.},label={lst:e-rows}]
always_ff @(posedge clk) begin
  sr_valid   <= {sr_valid[SR_ST-2:0], sr_feed};
  sr_rows[0] <= sr_row;
  sr_k[0]    <= sr_k_in;
  for (int s = 1; s < SR_ST; s++) begin
    sr_rows[s] <= sr_rows[s-1];
    sr_k[s]    <= sr_k[s-1];
  end
  for (int x = 0; x < 26; x++) begin
    sr_i1[x*5 +: 5] <= wirings_pkg::add26(5'(x), sr_k_in);
    sr_a[x*5 +: 5]  <= wirings_pkg::sub26(pick(tbl[0], sr_i1[x*5 +: 5]), sr_k[0]);
    sr_b[x*5 +: 5]  <= (sr_rows[1] >= 6'd26) ? pick(inner1, sr_a[x*5 +: 5])
                                             : pick(inner0, sr_a[x*5 +: 5]);
    sr_i4[x*5 +: 5] <= wirings_pkg::add26(sr_b[x*5 +: 5], sr_k[2]);
    sr_c[x*5 +: 5]  <= wirings_pkg::sub26(pick(tbl[1], sr_i4[x*5 +: 5]), sr_k[3]);
  end
end
\end{lstlisting}

다섯 단의 뜻은 식을 오른쪽에서 왼쪽으로 읽은 것이다. ① 글자 $x$에 $k$를 더한다. ② 오른쪽 로터의 정방향 표를
읽고 $k$를 뺀다(여기까지가 $R_k$). ③ 안쪽 치환을 읽는다($I$). ④ $k$를 더한다. ⑤ 오른쪽 로터의 역방향 표를
읽고 $k$를 뺀다($R_k^{-1}$). 앞 절과 같이 더하기와 ``읽고 빼기''를 다른 단에 두었다. \code{sr\_k}는 단마다
$k$를 한 칸씩 옮기는 꼬리표다. 다섯 단에 서로 다른 행이 들어 있으므로 각 단은 \emph{자기 행의} $k$를 써야
한다.

\Why 앞 절의 파이프라인은 글자를 한 줄로 흘렸는데 여기서는 \code{for (int x ...)}로 26줄을 나란히 놓았다.
같은 계산을 한 줄로 차례로 할지 여러 줄로 한꺼번에 할지는 \ref{sec:tradeoff}절의 맞바꿈이고, 답은 계산의
\emph{양}이 정한다. 안쪽 치환은 값이 26개뿐이어서 한 줄이면 26클럭이다. 행은 $26 \times 26 = 676$개의
값이다. 한 줄로 하면 676클럭이 들고, 26위치로 나누면 위치마다 26클럭이다. 전파 자체(위치마다 열 몇 클럭)보다
표 만드는 데 시간이 더 드는 셈이다. 26줄을 놓으면 26클럭에 끝나고 위치마다 1클럭이다.

\General 마트의 계산대와 같다. 손님이 적으면 계산대 하나로 충분하지만, 손님이 스물여섯 배가 되면 줄이
스물여섯 배 길어지거나 계산대를 스물여섯 개 열어야 한다. 회로의 예로는 오류 검출 부호(CRC)의 계산이 있다.
느린 직렬선에서는 비트 하나씩 처리하는 작은 회로를 쓰고, 빠른 통신망에서는 한 클럭에 32비트나 64비트를
한꺼번에 처리하는 넓은 회로를 쓴다. 계산은 같고 줄의 수가 다르다.

\Ours 그림~\ref{fig:g-e-lanes}가 이 구조다. 줄마다 표 조회가 세 번 있으므로 5비트 폭의 26:1 선택 회로가
$26 \times 3 = 78$개 든다. 엔진에서 표 계산이 차지하는 자원의 대부분이 여기다. 제어는 \code{INNER}와 같은
꼴이다. \code{sr\_row}를 0부터 25(또는 51)까지 올리며 한 클럭에 한 행씩 넣고, 단 수만큼 기다린다.
26번째부터의 행은 \code{sr\_rows[1] >= 26}이라는 조건으로 $I$ 대신 $I'$을 읽는다.

\begin{figure}[ht]
\centering
\begin{tikzpicture}[x=1cm,y=1cm]
  \foreach \lr/\lab in {0/{$x=$ A},1/{$x=$ B},2/{$x=$ C},4/{$x=$ Z}} {
    \node[anchor=east,font=\footnotesize] at (-0.1,-\lr*0.75) {\lab};
    \foreach \lc/\lt/\lst in {0/{$+k$}/blkb,1/{$R$}/blkg,2/{$I$}/blko,3/{$+k$}/blkb,4/{$R^{-1}$}/blkg} {
      \node[\lst,minimum width=13mm,minimum height=5.5mm,font=\footnotesize,inner sep=1pt] at (\lc*2.0+0.9,-\lr*0.75) {\lt};
      \draw[fill=cyellowbg] (\lc*2.0+1.65,-\lr*0.75-0.275) rectangle ++(0.16,0.55);
    }
    \foreach \lc in {0,1,2,3} { \draw[->] (\lc*2.0+1.81,-\lr*0.75) -- ++(0.42,0); }
  }
  \node at (4.9,-2.25) {$\vdots$};
  \draw[decorate,decoration={brace,amplitude=5pt}] (9.95,0.3) -- (9.95,-3.3);
  \node[reg,minimum width=26mm,minimum height=16mm] (ram) at (12.3,-1.5) {행 메모리\\{\scriptsize 130비트 $\times$ 64행}};
  \draw[->] (10.2,-1.5) -- node[above,note] {한 행} (ram);
  \node[note] at (4.9,-3.9) {노란 막대가 레지스터. 한 클럭마다 $k$가 하나씩 다른 행이 한 단씩 전진한다.};
\end{tikzpicture}
\caption{26레인의 행 계산. 가로가 파이프라인의 다섯 단, 세로가 26개의 레인이다. 끝에서 26개의 5비트 값이
모여 130비트 한 행이 되고 메모리의 $k$번 칸에 써진다.}
\label{fig:g-e-lanes}
\end{figure}

\section{선마다 하나씩: 분산 RAM의 사본}\label{sec:e-ram}

\begin{lstlisting}[language=SV,caption={\code{bombe\_engine.sv}: 메뉴의 선마다 행 메모리의 사본을 둔다.},label={lst:e-ram}]
for (genvar e = 0; e < MAX_EDGES; e++) begin : g_rows
  (* ram_style = "distributed" *) logic [129:0] rows [64];
  logic [5:0] rd_addr;
  assign rd_addr = {1'b0, wirings_pkg::add26(r, posm[e])} + (turn_e[e] ? 6'd26 : 6'd0);
  always_ff @(posedge clk) begin
    if (sr_valid[SR_ST-1]) rows[sr_rows[SR_ST-1]] <= sr_c;
    if (latch_perm)  perm_q[e*130 +: 130] <= rows[rd_addr];
  end
end
\end{lstlisting}

이 코드는 이상해 보인다. 방금 만든 행들을 메모리 하나에 넣으면 될 것을, 메뉴의 선 16개마다 \emph{똑같은
내용의} 메모리를 하나씩 두고 16개에 동시에 쓴다.

\Why 이유는 메모리가 레지스터와 다른 점에 있다. 레지스터는 값이 선으로 나와 있어 누구나 동시에 볼 수 있다.
메모리는 좁은 자리에 많은 값을 담는 대신, 값을 꺼내는 \term{단자}{port}의 수가 정해져 있다. 단자 하나는
한 번에 주소 하나의 값만 내놓는다. FPGA의 분산 RAM은 쓰는 단자 하나와 읽는 단자 하나를 가진 꼴이 기본이다.

그런데 메뉴의 선들은 같은 순간에 \emph{서로 다른} 행을 필요로 한다. 선 $e$가 크립의 $p$번째 글자에서 나온
것이면 그 선의 스크램블러는 오른쪽 로터가 $p$칸 더 간 위치에 있다(실제 봄베에서 유닛마다 드럼을 크립
위치만큼 어긋나게 끼우는 것에 해당한다). 그래서 선 $e$가 읽을 행은 $(r + p_e) \bmod 26$번이다. 코드의
\code{rd\_addr}이 이것이고, 넘어감 가설에서 넘어간 뒤쪽의 선이면 26을 더해 $I'$으로 만든 행을 읽는다.
선이 16개이면 주소도 16개다. 읽는 단자가 하나인 메모리로는 16클럭에 걸쳐 하나씩 읽을 수밖에 없다.

길은 셋이다(표~\ref{tab:g-e-ram}). 엔진은 셋째를 골랐다. 내용이 같은 메모리를 선마다 하나씩 두면 선마다
자기 읽는 단자가 생긴다. 쓰는 쪽은 모든 사본에 같은 값을 같은 주소로 쓰므로 회로가 늘지 않는다.

\begin{table}[ht]
\centering
\caption{16개의 선이 서로 다른 행을 읽게 하는 세 가지 방법.}
\label{tab:g-e-ram}
\small
\begin{tabularx}{\linewidth}{l>{\raggedright\arraybackslash}X>{\raggedright\arraybackslash}X}
\toprule
방법 & 드는 것 & 문제 \\
\midrule
메모리 하나를 차례로 읽는다 & 메모리 한 벌 & 위치마다 16클럭이 더 든다. 전파보다 읽기가 오래 걸린다 \\
행을 모두 레지스터에 둔다 & 플립플롭 $52 \times 130 = 6{,}760$개와, 선마다 130비트 폭의 52:1 선택 회로 &
  선택 회로가 $16 \times 130$벌이다(\ref{sec:mux-cost}절) \\
\textbf{메모리의 사본을 선마다 둔다} & 메모리 16벌 & 같은 내용을 16번 담는다. 그래도 메모리는 좁다 \\
\bottomrule
\end{tabularx}
\end{table}

\General 도서관의 사전 한 권을 열여섯 명이 동시에 서로 다른 쪽을 펴서 볼 수는 없다. 차례를 기다리거나
사전을 열여섯 권 사야 한다. 회로의 예로는 처리 장치의 레지스터 파일이 있다. 명령 하나가 레지스터 둘을
동시에 읽어야 하므로 읽는 단자가 둘 필요한데, FPGA에서는 흔히 같은 내용의 메모리 둘을 두고 쓸 때에는
양쪽에 같이 쓰는 것으로 이것을 만든다.

\Ours \code{(* ram\_style = "distributed" *)}는 합성 도구에게 이 배열을 분산 RAM으로 만들라고 알려 주는
표시다. 블록 RAM이 아니라 분산 RAM을 쓰는 이유는 둘이다. 64행뿐인 작은 메모리이고, 분산 RAM은 주소를 주면
같은 클럭 안에 값이 나온다(\ref{ch:resources}장). \code{genvar}로 쓴 \code{for}는 회로를 복사하라는 뜻이므로
(\ref{ch:hdl}장) 이 블록 전체가 선의 수만큼 생긴다. 합성 결과의 LUT RAM 2,858개는 주로 이 사본들이다.

\code{LATCH} 상태의 한 클럭 동안 모든 선이 자기 행을 읽어 \code{perm\_q}에 담는다. \code{perm\_q}는
선 16개 $\times$ 130비트의 레지스터로, \emph{지금 위치에서} 각 선의 스크램블러가 하는 치환이다. 기준 RTL에서
스크램블러 회로 16벌이 조합 논리로 내놓던 바로 그 값이 여기서는 메모리에서 읽은 값이다. 오른쪽 로터가 한
칸 갈 때마다 \code{LATCH}를 한 번 거치면 주소가 모두 하나씩 올라가 다음 행을 읽는다.

\section{케이블 단위 전파와 원핫 디코드}\label{sec:e-spread}

이제 전류를 퍼뜨린다. 계산의 내용은 기준 RTL의 \code{bombe\_core}와 한 글자도 다르지 않다
(\ref{ch:rtl-core}장). 다른 것은 회로의 구조다. \ref{ch:lesson}장에서 본 대로 첫 구현은 676비트를 한 줄로
놓고 $26a + y$라는 정수로 비트를 골랐고, 그 한 줄이 676:1 선택 회로와 676가닥 해독 회로를 832벌 만들어
LUT 40만 4천 개가 되었다. 고친 코드는 상태를 처음부터 2차원으로 둔다.

\begin{lstlisting}[language=SV,caption={\code{bombe\_engine.sv}: 케이블 단위의 전파(한 걸음의 다음 상태).},label={lst:e-spread}]
logic [25:0] cab_q [26];          // cable c, bit w = wire w live
logic [25:0] dec_a [MAX_EDGES];   // one-hot plain letter of edge i (0 if disabled)
logic [25:0] dec_b [MAX_EDGES];   // one-hot cipher letter of edge i (0 if disabled)

always_comb begin
  for (int i = 0; i < MAX_EDGES; i++) begin
    va[i] = cab_q[edge_q[i][4:0]];
    vb[i] = cab_q[edge_q[i][12:8]];
    for (int x = 0; x < 26; x++) begin
      to_b[i][x] = va[i][perm_q[(i*26 + x)*5 +: 5]];
      to_a[i][x] = vb[i][perm_q[(i*26 + x)*5 +: 5]];
    end
  end
  for (int c = 0; c < 26; c++) begin
    cab_n[c] = cab_q[c];
    for (int i = 0; i < MAX_EDGES; i++) begin
      if (dec_b[i][c]) cab_n[c] = cab_n[c] | to_b[i];
      if (dec_a[i][c]) cab_n[c] = cab_n[c] | to_a[i];
    end
    if (diagonal_en) begin
      for (int w = 0; w < 26; w++) cab_n[c][w] = cab_n[c][w] | cab_q[w][c];
    end
  end
end
\end{lstlisting}

한 줄씩 읽는다.

\begin{itemize}
\item \textbf{\code{cab\_q}.} 26비트 낱말 26개다. \code{cab\_q[c][w]}가 케이블 $c$의 $w$번 선이다. 합쳐서
676비트인 것은 같지만 ``어느 케이블''과 ``어느 선''이 따로 있다.
\item \textbf{\code{va[i] = cab\_q[...]}.} 선 $i$의 한쪽 글자에 해당하는 케이블을 \emph{통째로} 꺼낸다.
26개의 낱말 가운데 하나를 고르는 26비트 폭의 26:1 선택이다. \code{vb[i]}는 반대쪽 글자의 케이블이다.
\item \textbf{\code{to\_b[i][x] = va[i][perm...]}.} 꺼낸 26가닥을 스크램블러의 치환으로 자리를 옮긴다.
건너편 케이블의 $x$번 선에 전류를 보내는 것은 이쪽의 $S(x)$번 선이다. 에니그마의 치환은 거꾸로 해도 같은
치환이어서($S(S(x)) = x$) ``$x$로 가는 선''과 ``$x$에서 가는 선''이 같은 번호다. 그래서 표 하나로 양쪽
방향을 다 쓴다. 출력의 선마다 26가닥 가운데 하나를 고르는 26:1 선택이다.
\item \textbf{\code{if (dec\_b[i][c]) ...}.} 자리를 옮긴 26가닥을 목적지 케이블에 OR로 더한다. 케이블 $c$의
처지에서 보면 ``선 $i$의 건너편이 나인가''를 묻고, 그렇다면 그 26가닥을 받는다.
\item \textbf{대각선 보드.} \code{cab\_n[c][w] | cab\_q[w][c]}. 케이블 $c$의 $w$번 선과 케이블 $w$의 $c$번
선을 잇는다. 번호가 상수이므로 선택 회로가 없다. 칩 안에서는 전선 한 가닥과 OR의 입력 하나일 뿐이다.
실제 봄베에서 대각선 보드가 릴레이도 스위치도 없는 고정 배선이었던 것과 같다.
\end{itemize}

\paragraph{원핫 디코드.} \code{dec\_b[i]}는 26비트 가운데 한 비트만 1인 값이다. 선 $i$의 건너편 글자가
E이면 E 자리의 비트만 1이다. 이런 표현을 \term{원핫}{one-hot}이라 한다.

\Why ``선 $i$의 목적지가 케이블 $c$인가''를 판단하는 방법은 둘이다. 글자 번호(5비트)를 $c$와 비교하는
회로를 그 자리에 두거나, 미리 26가닥으로 풀어 둔 신호에서 $c$번 가닥을 보거나. 앞쪽은 비교 회로를 지나야
답이 나오고, 뒤쪽은 이미 답이 선 위에 있다. 전파 회로는 상태 레지스터에서 나와 상태 레지스터로 돌아가는
고리이고 이 고리의 길이가 엔진의 클럭을 정한다(\ref{sec:timing}절). 고리 안에서 할 일을 고리 밖으로 꺼낼 수
있으면 꺼내야 한다. 메뉴는 구동 중에 바뀌지 않으므로 이 풀이는 구동을 시작할 때 한 번 해서 레지스터에
넣어 둔다.

\begin{lstlisting}[language=SV,numbers=none]
dec_a[e] <= edge_en_q[e] ? (26'd1 << edge_q[e][4:0])  : 26'd0;
dec_b[e] <= edge_en_q[e] ? (26'd1 << edge_q[e][12:8]) : 26'd0;
\end{lstlisting}

쓰지 않는 선은 26비트가 모두 0이 된다. 어느 케이블에도 전류를 보내지 않으므로 ``이 선이 켜져 있는가''를
따로 검사할 필요가 없다.

\General 우체국에서 편지마다 주소를 읽어 가며 칸에 넣는 것과, 미리 칸 번호 스티커를 붙여 둔 편지를 칸마다
``내 번호인가''만 보고 집는 것의 차이다. 회로의 예로는 메모리 여러 개를 쓰는 시스템의 칩 선택 신호가 있다.
주소의 윗자리를 한 번 풀어 메모리마다 선 하나씩을 주면, 각 메모리는 자기 선 하나만 본다. 상태 기계의 상태를
원핫으로 두는 것도 같은 생각이다. ``지금 상태가 3번인가''가 비교가 아니라 선 하나가 된다.

\Ours 그림~\ref{fig:g-e-cable}은 케이블 하나가 다음 값을 얻는 회로다. 선마다 양쪽 끝이 있으므로 받을 수
있는 묶음은 최대 32개이고, 케이블은 그 가운데 자기 앞으로 온 것만 OR로 모은다. 이 구조로 엔진 전체가
LUT 41,040개가 되었다. 첫 구현의 약 10분의 1이다(그림~\ref{fig:g-lutbars}).

\begin{figure}[ht]
\centering
\begin{tikzpicture}[x=1cm,y=1cm]
  \foreach \i/\yy/\gt/\gd in {0/2.4/{\code{to\_b[0]}}/{\code{dec\_b[0][c]}},1/1.2/{\code{to\_a[0]}}/{\code{dec\_a[0][c]}},3/-0.6/{\code{to\_a[15]}}/{\code{dec\_a[15][c]}}} {
    \node[blkb,minimum width=24mm,minimum height=7mm,font=\footnotesize] (t\i) at (0,\yy) {\gt};
    \node[and gate US,draw,fill=white,logic gate inputs=nn,minimum height=7mm] (g\i) at (4.2,\yy) {};
    \draw[->,line width=1.4pt,cgray] (t\i.east) -- node[above,font=\scriptsize,text=cgray,pos=0.3] {26} (g\i.input 1);
    \node[font=\scriptsize,anchor=east] (d\i) at (3.1,\yy-0.45) {\gd};
    \draw[->] (d\i.east) -| ([xshift=-2mm]g\i.input 2) -- (g\i.input 2);
  }
  \node at (0,0.4) {$\vdots$}; \node at (4.2,0.4) {$\vdots$};
  \node[or gate US,draw,fill=white,logic gate inputs=nnnnn,minimum height=40mm] (o) at (8.2,-0.2) {};
  \draw[->,line width=1.4pt,cgray] (g0.output) -- (5.6,2.4) |- (o.input 1);
  \draw[->,line width=1.4pt,cgray] (g1.output) -- (5.3,1.2) |- (o.input 2);
  \draw[->,line width=1.4pt,cgray] (g3.output) -- (5.3,-0.6) |- (o.input 3);
  \node[blkn,font=\footnotesize,minimum width=40mm] (dg) at (2.6,-1.8) {대각선 보드: \code{cab\_q[w][c]}};
  \node[blkn,font=\footnotesize,minimum width=40mm] (old) at (2.6,-2.8) {지금의 값: \code{cab\_q[c]}};
  \draw[->,line width=1.4pt,cgray] (dg.east) -- (5.9,-1.8) |- (o.input 4);
  \draw[->,line width=1.4pt,cgray] (old.east) -- (6.2,-2.8) |- (o.input 5);
  \node[reg,minimum width=22mm,minimum height=12mm] (n) at (11.6,-0.2) {\code{cab\_n[c]}\\{\scriptsize 26비트}};
  \draw[->,line width=1.4pt,cgray] (o.output) -- (n);
  \node[note,anchor=west,text width=34mm,align=left] at (9.9,-2.4) {이 그림이 케이블마다 한 벌, 모두 26벌};
\end{tikzpicture}
\caption{케이블 $c$의 다음 값. 굵은 선은 26가닥 묶음이다. 원핫 신호 한 가닥이 묶음 전체를 통과시키거나
막는다(AND). 통과한 묶음과 지금 값과 대각선 보드에서 온 값을 OR로 모은다.}
\label{fig:g-e-cable}
\end{figure}

\section{고정점 판정은 레지스터에서}\label{sec:e-fixed}

전파는 더 켜질 선이 없을 때까지, 곧 한 걸음을 더 가도 상태가 바뀌지 않을 때까지 되풀이한다. 그 상태가
\term{고정점}{fixed point}이다. ``바뀌지 않았다''를 어떻게 알아내는가.

\begin{lstlisting}[language=SV,caption={\code{bombe\_engine.sv}: 판정에 쓰는 신호와 \code{SPREAD} 상태.},label={lst:e-fixed}]
assign test_q    = cab_q[test_letter];
assign test_full = (test_q == FULL);
always_comb begin
  stable = 1'b1;
  for (int c = 0; c < 26; c++) stable = stable & (cab_q[c] == cab_prev[c]);
end

S_SPREAD: begin
  if (stable || test_full) begin
    positions_q <= positions_q + 32'd1;
    state <= test_full ? S_ADVANCE : S_PUSH;
  end else begin
    cab_prev <= cab_q;
    cab_q    <= cab_n;
  end
end
\end{lstlisting}

\Why 가장 먼저 떠오르는 방법은 다음 값과 지금 값을 비교하는 것이다(\code{cab\_n == cab\_q}). 그런데
\code{cab\_n}은 앞 절의 전파 회로 전체를 지나서 나오는 값이고, 그 회로가 이미 엔진에서 가장 긴 경로다.
그 끝에 676비트 비교 회로를 더 붙이고 그 결과로 상태 기계까지 움직이면 가장 긴 경로가 더 길어진다.
엔진은 대신 \emph{이미 레지스터에 들어 있는 두 값}을 비교한다. 한 걸음 전의 상태를 \code{cab\_prev}에
따로 남겨 두고, \code{cab\_q}와 \code{cab\_prev}가 같은지를 본다. 두 값 모두 레지스터에서 바로 나오므로
이 비교는 클럭 주기의 처음부터 시작할 수 있고 전파 회로와 겹치지 않는다. 대가는 한 클럭이다. 상태가 더
바뀌지 않게 된 것을 한 클럭 뒤에야 안다(그림~\ref{fig:g-e-fixed}).

\General 연못의 물결이 잦아들었는지 알려면 지금 찍은 사진과 조금 전에 찍은 사진을 나란히 놓고 보면 된다.
아직 찍지 않은 사진과 비교하려고 기다릴 필요가 없다. 회로의 예로는 계수기의 끝 검출이 있다. ``다음 값이
끝값인가''를 덧셈기의 출력에서 검사하면 덧셈기 뒤에 비교 회로가 이어진다. ``지금 값이 끝값인가''를
레지스터에서 검사하면 덧셈기와 비교 회로가 나란히 놓인다. 한 클럭 늦게 아는 대신 경로가 짧다.

\begin{figure}[ht]
\centering
\begin{tikztimingtable}[timing/slope=0.1,timing/coldist=4pt,xscale=2.2,yscale=1.25,semithick]
  \code{clk}       & 8{.5L .5H} \\
  상태             & .5D{LATCH} 5D{SPREAD} 1D{PUSH} 1D{ADV.} .5D{LATCH} \\
  \code{cab\_q}    & .5U 1D{$L_0$} 1D{$L_1$} 1D{$L_2$} 4D{$L_3$} .5U \\
  \code{cab\_prev} & .5U 1D{전부 1} 1D{$L_0$} 1D{$L_1$} 1D{$L_2$} 3D{$L_3$} .5U \\
  \code{stable}    & 4.5L 3H .5L \\
\extracode
  \begin{pgfonlayer}{background}
    \vertlines[cgray!50,dashed]{0.5,1.5,...,7.5}
    \fill[cgreenbg] (4.5,-9) rectangle (5.5,1.2);
  \end{pgfonlayer}
  \node[note,text=cgreen] at (5.0,-9.9) {여기서 판정};
\end{tikztimingtable}
\caption{한 위치의 진행(정지인 경우). $L_0$은 시험 선 하나만 켜진 처음 상태다. $L_3$에서 더 변하지 않는데,
\code{cab\_q}와 \code{cab\_prev}가 같아지는 것은 그다음 클럭이다. 실제로는 걸음 수가 더 많다.}
\label{fig:g-e-fixed}
\end{figure}

\Ours 세 가지를 짚는다.

\begin{itemize}
\item \textbf{처음 값.} \code{LATCH}에서 \code{cab\_q}는 시험 선 하나만 켜고 \code{cab\_prev}는 전부 1로
채운다. 전부 1인 상태는 통과로 끝나는 경우에만 나올 수 있으므로, 첫 클럭에 ``같다''로 잘못 판정되는 일이
없다. 앞 위치의 \code{cab\_prev}가 우연히 지금의 처음 값과 같을 가능성을 없애는 것이다.
\item \textbf{조기 종료.} \code{test\_full}은 테스트 글자의 케이블 26선이 전부 켜졌다는 뜻이다. 그러면
가설은 모순이고 이 위치는 통과다. 켜진 선은 다시 꺼지지 않으므로 고정점까지 기다릴 필요가 없다. 대부분의
위치가 여기서 끝난다. 기준 RTL과 Rust 코어도 같은 조기 종료를 쓴다.
\item \textbf{정지의 조건.} 고정점에 이르렀는데 테스트 레지스터가 가득 차지 않았으면 정지다. 코드에서
\code{test\_full}이면 \code{ADVANCE}로, 아니면 \code{PUSH}로 간다. 실제 봄베의 감지 릴레이 26개가 하던
일이 이 한 줄이다.
\end{itemize}

한 위치에 드는 클럭은 \code{LATCH} 1, 전파 걸음 수, 판정 1, \code{ADVANCE} 1이고 여기에 표 계산을 나눈
몫 약 3이 더해진다. 시뮬레이션에서 잰 평균은 메뉴에 따라 위치당 12에서 17클럭이다.

\section{정지 FIFO}\label{sec:e-fifo}

정지가 나오면 드럼 위치 네 자리와 테스트 레지스터의 값을 호스트에 넘겨야 한다. 호스트는 느리다. 직렬선으로
값 하나를 읽어 가는 데 엔진의 클럭으로 수천 번이 지난다. 그동안 엔진을 세워 두지 않으려고 FIFO를 둔다
(\ref{sec:handshake}절).

\begin{lstlisting}[language=SV,caption={\code{bombe\_engine.sv}: 정지 FIFO.},label={lst:e-fifo}]
localparam int DEPTH = 1 << FIFO_LOG2;        // 64 entries
logic [45:0]        fifo [DEPTH];             // {g, l, m, r, test register}
logic [FIFO_LOG2:0] wr_ptr, rd_ptr;           // one bit wider than the address
assign fifo_empty = (wr_ptr == rd_ptr);
assign fifo_full  = (wr_ptr[FIFO_LOG2] != rd_ptr[FIFO_LOG2]) &&
                    (wr_ptr[FIFO_LOG2-1:0] == rd_ptr[FIFO_LOG2-1:0]);

always_ff @(posedge clk) begin
  if (push) fifo[wr_ptr[FIFO_LOG2-1:0]] <= {g, l, m, r, test_q};
end
\end{lstlisting}

칸 하나는 46비트다. 드럼 위치 네 자리가 5비트씩 20비트, 테스트 레지스터가 26비트다. 칸은 64개다.

\Why 그림~\ref{fig:g-fifo}에서 ``두 위치가 같으면 비어 있다''고 했다. 그런데 64칸을 다 채워도 쓰기 위치가
한 바퀴 돌아 읽기 위치와 같아진다. 비어 있는 것과 가득 찬 것을 구별해야 한다. 이 코드의 방법은 위치
계수기를 주소보다 \emph{한 비트 넓게} 두는 것이다. 아래 여섯 비트가 칸의 주소이고 맨 윗비트는 ``몇 바퀴째인가''의
홀짝이다. 아래 비트가 같을 때 윗비트까지 같으면 같은 바퀴이므로 비어 있고, 윗비트가 다르면 쓰기 쪽이 한
바퀴 앞선 것이므로 가득 찼다.

\General 운동장 트랙을 도는 두 사람이 같은 지점에 있을 때, 나란히 달리는 중인지 한 사람이 한 바퀴 앞섰는지는
자리만 보아서는 모른다. 바퀴 수를 세고 있으면 안다. 여기서는 바퀴 수의 홀짝만 세면 충분하다. 한 바퀴보다
더 앞설 수는 없기 때문이다.

\Ours 가득 찼을 때 엔진은 정지를 버리지 않고 \code{PUSH} 상태에서 기다린다. 코드~\ref{lst:e-fsm-push}의
\code{if (!fifo\_full)}이 그것이다. 호스트가 하나를 읽어 가면 자리가 나고 그때 넣고 다음 위치로 간다.
\ref{sec:handshake}절의 역압이다. 그래서 정지 목록은 호스트가 아무리 느려도 빠짐없이, 순서대로 전달된다.

\begin{lstlisting}[language=SV,caption={\code{bombe\_engine.sv}: \code{PUSH}와 \code{ADVANCE}.},label={lst:e-fsm-push}]
S_PUSH: begin
  if (!fifo_full) begin
    wr_ptr <= wr_ptr + 1'b1;
    state  <= S_ADVANCE;
  end
end
S_ADVANCE: begin
  if (last) begin
    done_q <= 1'b1;
    state  <= S_IDLE;
  end else if (r != 5'd25) begin
    r     <= r + 5'd1;
    state <= S_LATCH;          // only the right rotor moved: next row
  end else begin
    r <= '0;
    // ... carry into m, l, g as in bombe_sweep ...
    in_x     <= '0;
    wait_cnt <= '0;
    state    <= S_INNER;       // the middle rotor moved: rebuild the tables
  end
end
\end{lstlisting}

\code{ADVANCE}는 \ref{ch:rtl-sweep}장의 26진 자리올림과 같은데, 어디로 돌아가는지가 두 갈래다. 오른쪽
로터만 한 칸 갔으면 표는 그대로이므로 \code{LATCH}로 가서 다음 행을 읽는다. 자리올림이 일어나 가운데
로터(또는 그 위)가 바뀌었으면 \code{INNER}로 가서 표를 다시 만든다. \ref{sec:e-recompute}절의 층 구조가
코드에서는 이 두 갈래다. 새 구동을 시작할 때에는 \code{wr\_ptr <= rd\_ptr}로 쓰기 위치를 읽기 위치에 맞춰,
앞 구동에서 읽어 가지 않은 정지를 버린다.

\chapter{엔진 여러 대와 바깥 세계: \texorpdfstring{\code{bombe\_array}}{bombe\_array}, 레지스터 버스, 검증}\label{ch:fpga-array}

엔진은 혼자서는 아무것도 하지 못한다. 누군가 메뉴와 로터 순서를 써넣고, 시작시키고, 정지를 읽어 가야 한다.
이 장은 엔진과 바깥을 잇는 부분과, 엔진이 옳다는 근거와, 칩에 올린 결과를 다룬다.

\section{레지스터 버스}\label{sec:e-bus}

\Why 회로의 안쪽에는 설정값과 상태값이 레지스터에 담겨 있다. 바깥의 프로그램이 이것을 읽고 쓰게 하는 가장
흔한 방법은 레지스터마다 \emph{주소}를 붙이는 것이다. 바깥에서는 ``몇 번 주소에 이 값을 써라'', ``몇 번
주소의 값을 읽어 달라''는 두 가지 요청만 한다. 회로 쪽은 주소를 보고 해당하는 레지스터에 값을 넣거나
꺼내 준다. 이렇게 하면 회로에 기능이 늘어도 연결선은 그대로이고 주소만 늘어난다. 이런 연결을
\term{버스}{bus}, 주소가 붙은 레지스터들의 목록을 \term{레지스터 맵}{register map}이라 한다.

\General 아파트의 우편함이 주소가 붙은 레지스터다. 집배원은 집 안의 사정을 몰라도 호수만 보고 넣는다.
회로의 예로는 마이크로컨트롤러의 주변 장치가 모두 이 방식이다. 프로그램이 특정 주소에 1을 쓰면 LED가
켜지고, 다른 주소를 읽으면 버튼의 상태가 나온다. 프로그램에게는 메모리를 읽고 쓰는 것과 똑같아 보인다.

\Ours 엔진의 버스는 신호가 여섯이다. 요청(\code{bus\_req}), 쓰기인지 읽기인지(\code{bus\_we}),
주소(\code{bus\_addr}), 쓸 값(\code{bus\_wdata})이 들어가고, 응답(\code{bus\_ack})과 읽은
값(\code{bus\_rdata})이 나온다. 요청은 한 클럭 동안 올리고 응답은 그다음 클럭에 온다
(그림~\ref{fig:g-e-bus}). 레지스터 맵은 표~\ref{tab:g-e-regmap}과 같다.

\begin{figure}[ht]
\centering
\begin{tikztimingtable}[timing/slope=0.1,timing/coldist=4pt,xscale=2.4,yscale=1.25,semithick]
  \code{clk}        & 7{.5L .5H} \\
  \code{bus\_req}   & .5L 1H 2L 1H 2.5L \\
  \code{bus\_we}    & .5L 1H 5.5L \\
  \code{bus\_addr}  & .5U 1D{0x008} 2U 1D{0x004} 2.5U \\
  \code{bus\_wdata} & .5U 1D{값} 5.5U \\
  \code{bus\_ack}   & 1.5L 1H 2L 1H 1.5L \\
  \code{bus\_rdata} & 4.5U 1D{상태} 1.5U \\
\extracode
  \begin{pgfonlayer}{background}
    \vertlines[cgray!50,dashed]{0.5,1.5,...,6.5}
  \end{pgfonlayer}
  \node[note] at (1.5,-13.6) {쓰기};
  \node[note] at (4.5,-13.6) {읽기};
\end{tikztimingtable}
\caption{레지스터 버스의 쓰기 한 번과 읽기 한 번. 응답은 요청의 다음 클럭에 온다.}
\label{fig:g-e-bus}
\end{figure}

\begin{table}[ht]
\centering
\caption{엔진의 레지스터 맵(\code{rtl/fpga/README.md}의 요약).}
\label{tab:g-e-regmap}
\small
\begin{tabularx}{\linewidth}{llc>{\raggedright\arraybackslash}X}
\toprule
주소 & 이름 & 방향 & 내용 \\
\midrule
\code{0x000} & \code{CTRL} & 쓰기 & 시작, 중단 \\
\code{0x004} & \code{STATUS} & 읽기 & 구동 중, 끝남, 정지가 FIFO에 있음 \\
\code{0x008} & \code{ORDER} & 읽기/쓰기 & 반사판과 로터 넷의 번호 \\
\code{0x00C} & \code{OPTIONS} & 읽기/쓰기 & 테스트 글자, 시험 선, 대각선 보드, 넘어감 가설 \\
\code{0x010} & \code{EDGE\_EN} & 읽기/쓰기 & 메뉴의 선 가운데 쓰는 것 \\
\code{0x014}, \code{0x018} & \code{POSITIONS}, \code{CYCLES} & 읽기 & 시험한 위치 수, 지난 클럭 수 \\
\code{0x01C} & \code{STOP\_POS} & 읽기 & FIFO 맨 앞 정지의 드럼 위치 \\
\code{0x020} & \code{STOP\_REG} & 읽기 & 그 정지의 테스트 레지스터. \emph{읽으면 FIFO에서 빠진다} \\
\code{0x100}부터 & \code{EDGE[i]} & 읽기/쓰기 & 선 $i$의 두 글자와 크립 위치 \\
\bottomrule
\end{tabularx}
\end{table}

\begin{lstlisting}[language=SV,caption={\code{bombe\_engine.sv}: 버스의 쓰기와 읽기(일부).},label={lst:e-bus}]
bus_ack     <= bus_req;
start_pulse <= 1'b0;
abort_pulse <= 1'b0;
if (bus_req && bus_we) begin
  case (bus_addr)
    12'h000: begin
      start_pulse <= bus_wdata[0];
      abort_pulse <= bus_wdata[1];
    end
    12'h008: order_q   <= bus_wdata[19:0];
    // ...
  endcase
end
if (bus_req && !bus_we) begin
  case (bus_addr)
    12'h004: bus_rdata <= {28'd0, overflow_q, !fifo_empty, done_q, busy};
    12'h020: begin
      bus_rdata <= fifo_empty ? 32'd0 : {6'd0, fifo[rd_ptr[FIFO_LOG2-1:0]][25:0]};
      if (!fifo_empty) rd_ptr <= rd_ptr + 1'b1;
    end
    // ...
  endcase
end
\end{lstlisting}

코드에서 볼 것이 셋이다.

\begin{itemize}
\item \textbf{응답.} \code{bus\_ack <= bus\_req}. 요청을 한 클럭 늦춘 것이 응답이다. 읽은 값도 같은
클럭에 레지스터에 담기므로 응답과 값이 함께 나온다.
\item \textbf{펄스.} \code{CTRL}은 값을 저장하는 레지스터가 아니다. \code{start\_pulse}는 매 클럭 0으로
돌아가도록 먼저 적고, 쓰기가 왔을 때에만 그 클럭에 1이 되게 덮어쓴다. 같은 블록에서 뒤에 적은 대입이
이긴다는 규칙을 쓴 것이다. 결과는 한 클럭짜리 신호이고 상태 기계가 이것을 보고 \code{IDLE}을 떠난다.
``시작''은 상태가 아니라 사건이므로 이렇게 만든다.
\item \textbf{읽으면 바뀌는 레지스터.} \code{STOP\_REG}를 읽으면 읽기 위치가 하나 올라간다. 호스트는
\code{STOP\_POS}를 먼저 읽고 \code{STOP\_REG}를 읽어 정지 하나를 꺼낸다. 소프트웨어에서 ``읽기''는 보통
아무것도 바꾸지 않지만, 회로의 레지스터 맵에서는 읽기에 동작을 붙이는 일이 흔하다.
\end{itemize}

\section{엔진 여러 대: \texorpdfstring{\code{bombe\_array}}{bombe\_array}}\label{sec:e-array}

\begin{lstlisting}[language=SV,caption={\code{bombe\_array.sv}: 주소의 윗자리로 엔진을 고른다.},label={lst:e-array}]
assign sel = bus_addr[15:12];

for (genvar k = 0; k < N_ENGINES; k++) begin : g_engine
  bombe_engine #(.MAX_EDGES(MAX_EDGES), .FIFO_LOG2(FIFO_LOG2)) u_engine (
    .clk      (clk),
    .rst_n    (rst_n),
    .bus_req  (bus_req && sel == 4'(k)),
    .bus_addr (bus_addr[11:0]),
    // ...
    .bus_ack  (e_ack[k]),
    .bus_rdata(e_rdata[k])
  );
end

always_comb begin
  bus_ack   = g_ack | (|e_ack);
  bus_rdata = g_ack ? g_rdata : 32'd0;
  for (int k = 0; k < N_ENGINES; k++) begin
    if (e_ack[k]) bus_rdata = e_rdata[k];
  end
end
\end{lstlisting}

\Why 로터 순서는 60가지(Enigma~I)에서 1,344가지(M4)다. 로터 순서들은 서로 무관하므로 엔진을 여러 대
놓고 하나씩 맡기면 그만큼 빨리 끝난다. 프로그램에서 스레드를 여러 개 돌리는 것에 해당하지만, 회로에서는
같은 회로를 여러 벌 놓는 것이고 칩의 자리를 그만큼 차지한다.

\General 아파트 주소가 ``동''과 ``호''로 나뉘듯, 버스의 주소도 윗자리와 아랫자리로 나눈다. 윗자리로 장치를
고르고 아랫자리는 그 장치 안의 레지스터 번호로 넘긴다. 컴퓨터의 주소 공간이 메모리, 화면, 통신 장치에
구역별로 나뉘어 있는 것이 이 방식이다.

\Ours 16비트 주소의 위 네 비트가 엔진의 번호다. \code{genvar} 반복문이 엔진을 $N$대 만들고, 엔진 $k$에게는
번호가 $k$일 때에만 요청을 전한다. 아래 열두 비트는 그대로 넘기므로 엔진은 자기가 몇 번째인지 모른다.
모든 엔진이 같은 버스 선을 듣고 있고, 요청 신호만 한 대에게 간다. 응답은 한 대만 하므로 OR로 모으면 된다.
번호 15(\code{0xF000})는 엔진이 아니라 식별값과 엔진 수를 알려 주는 구역이고, 없는 번호의 엔진에 대한
요청에도 0으로 응답해 호스트가 기다리다 멈추지 않게 한다.

블레츨리 파크의 막사에 봄베 여러 대가 나란히 놓여 저마다 다른 로터 순서를 돌리던 모습과 같다. 호스트의
프로그램(\code{rtl/fpga/host})이 그 막사의 감독이다. 로터 순서의 목록을 엔진들에 나눠 주고, 끝난 엔진에
다음 것을 주고, 정지를 모아 Rust 코어의 체킹 단계로 넘긴다.

\section{브리지: 버스를 바깥으로}\label{sec:e-bridge}

레지스터 버스는 칩 안의 약속이다. 칩 밖의 컴퓨터와 이으려면 보드에 있는 통로에 맞는 번역기가 필요하다.
보드가 둘이므로 번역기도 둘이고, 엔진 쪽은 어느 것이 붙었는지 모른다.

\begin{itemize}
\item \textbf{\code{uart\_bridge}} (Nexys Video). PC와 직렬선(UART)으로 통신한다. 직렬선은 선 하나로 비트를
하나씩 보내는 가장 단순한 통로다. 브리지는 바이트를 받아 모아 버스의 요청 하나로 바꾼다. 쓰기는
\code{'W'}, 주소 2바이트, 값 4바이트의 일곱 바이트이고 브리지가 한 바이트로 답한다. 읽기는 \code{'R'}과
주소 2바이트이고 브리지가 값 4바이트로 답한다. 속도는 초당 921,600비트다. 바이트 사이가 너무 벌어지면 받던
명령을 버리고 처음으로 돌아가, 선에서 바이트 하나를 잃어도 그 뒤가 계속 어긋나지 않게 한다.
\item \textbf{\code{axil\_bridge}} (PYNQ). 이 보드의 칩에는 FPGA와 ARM 처리 장치가 함께 들어 있고, 둘
사이의 표준 버스가 AXI4-Lite다. 브리지는 AXI4-Lite의 쓰기와 읽기를 엔진의 버스 요청으로 바꾼다. 처리 장치의
프로그램에게 엔진의 레지스터는 메모리의 한 구역으로 보인다.
\end{itemize}

두 브리지의 코드는 \ref{sec:fsm}절의 상태 기계와 \ref{sec:handshake}절의 핸드셰이크를 그대로 쓴 것이어서
이 문서에서는 한 줄씩 따라가지 않는다.

\section{엔진이 옳다는 근거}\label{sec:e-verify}

\Why 엔진은 기준 RTL과 구조가 전혀 다르다. 표를 미리 만들고, 메모리에서 읽고, 케이블 단위로 전파하고,
한 클럭 늦게 판정한다. 어느 한 곳이라도 틀리면(꼬리표가 한 단 어긋나거나, 읽는 주소가 한 칸 밀리거나)
결과가 달라진다. 구조가 다를수록 ``같은 답을 낸다''는 것을 따로 확인해야 하고, 그 방법은
\ref{ch:verify}장에서 본 대로 기준 모델과의 대조다.

\Ours 시험(\code{rtl/fpga/tb/cocotb\_engine.py})은 엔진의 안쪽 신호를 건드리지 않는다. 실제 호스트가 하는
것과 똑같이 \emph{레지스터 버스로만} 엔진을 다룬다. 메뉴와 로터 순서를 레지스터에 쓰고, 시작시키고,
정지를 읽어 모은 뒤 Rust 코어가 내놓은 정지 목록과 비교한다(코드~\ref{lst:e-tb}). 비교 대상인 목록은
기준 RTL의 \code{bombe\_sweep}을 시험할 때 쓴 바로 그 파일이다. 그러므로 Rust 코어, 기준 RTL, FPGA 엔진의
셋이 같은 목록으로 묶인다(그림~\ref{fig:g-e-verify}).

\begin{lstlisting}[language=Python,caption={엔진 시험의 꼴(요약. 실제 코드는 \code{rtl/fpga/tb/cocotb\_engine.py}).},label={lst:e-tb}]
await configure(bus, base, config_of(ref))      # ORDER, OPTIONS, EDGE[i], EDGE_EN
await bus.write(base + CTRL, 1)                 # START
stops = []
while (await bus.read(base + STATUS)) & BUSY or stop_available:
    pos = await bus.read(base + STOP_POS)
    reg = await bus.read(base + STOP_REG)       # pops the FIFO
    stops.append((pos, reg))
assert stops == ref["stops"]                    # same stops, same order
assert await bus.read(base + POSITIONS) == ref["positions_tested"]
\end{lstlisting}

\begin{figure}[ht]
\centering
\begin{tikzpicture}[x=1cm,y=1cm]
  \node[blkg,minimum width=30mm,minimum height=12mm] (rust) at (0,0) {Rust 코어\\{\footnotesize \code{Bombe.run}}};
  \node[blkn,minimum width=34mm,minimum height=12mm] (vec) at (4.6,0) {정지 목록 파일\\{\footnotesize \code{data/vectors/rtl}}};
  \node[blkb,minimum width=34mm,minimum height=11mm] (ref) at (10.0,0.9) {기준 RTL \code{bombe\_sweep}};
  \node[blkb,minimum width=34mm,minimum height=11mm] (eng) at (10.0,-0.9) {FPGA 엔진 (버스로 구동)};
  \draw[->] (rust) -- (vec);
  \draw[<->,dashed] (vec.east) -- ++(0.6,0) |- node[above,note,pos=0.75] {순서까지 비교} (ref.west);
  \draw[<->,dashed] (vec.east) -- ++(0.6,0) |- node[below,note,pos=0.75] {순서까지 비교} (eng.west);
  \node[blko,minimum width=34mm,minimum height=11mm] (uart) at (10.0,-2.6) {UART 브리지 + 엔진 2대};
  \draw[<->,dashed] (eng) -- node[right,note] {같은 시험을 직렬선으로} (uart);
\end{tikzpicture}
\caption{검증의 연결. 같은 파일이 기준 RTL과 FPGA 엔진의 답안지다.}
\label{fig:g-e-verify}
\end{figure}

시험은 다섯 항목이고 모두 통과한다.

\begin{itemize}
\item \textbf{식별 구역과 레지스터.} \code{0xF000}의 식별값, 엔진 수, 쓴 값이 그대로 읽히는지.
\item \textbf{Enigma~I 전체 스윕 두 묶음.} 엔진 두 대에 서로 다른 메뉴를 걸고 \emph{동시에} 돌린다.
각각 17,576 위치를 훑어 정지 목록이 Rust 코어와 순서까지 같고 시험한 위치 수도 같은지 본다. 묶음 가운데
하나는 대각선 보드를 끄고 넘어감 가설을 켠 메뉴다. 두 대를 함께 돌리는 것은 한 엔진의 요청이 다른 엔진에
새지 않는지를 보기 위해서다.
\item \textbf{M4.} 그리스 로터가 있는 4로터 기종에서 앞의 정지 60개가 순서대로 같은지 보고 중단시킨다.
전체 456,976 위치는 시뮬레이션으로는 오래 걸리므로 선택 사항이다.
\item \textbf{중단과 재시작.} 중단한 뒤 다시 시작하면 앞 구동의 남은 정지가 섞이지 않고 처음부터 같은
정지가 나오는지.
\end{itemize}

정지 목록을 \emph{순서까지} 비교한다는 점이 중요하다. 정지의 집합만 같으면 되는 것 아닌가 싶지만, 순서가
같다는 것은 드럼 위치를 넘기는 순서, 표를 다시 만드는 시점, FIFO가 모두 맞다는 뜻이다. 위치 수의 비교는
자리올림이 틀려 일부 위치를 건너뛴 경우를 잡는다(\ref{ch:rtl-sweep}장).

브리지도 따로 시험한다. UART 브리지에 엔진 두 대를 붙이고 직렬선의 바이트만으로 같은 구동을 시키는 종단
시험이 있고, PC 쪽 호스트 프로그램의 시험 열두 개가 통과한다. 여기까지가 시뮬레이션이다. 실제 보드에서의
시험은 보드를 확보한 뒤의 일이다.

\section{칩에 올린 결과}\label{sec:e-results}

\begin{lesson}{클럭을 절반으로}
케이블 단위 구조로 고친 뒤 합성 단계의 추정은 100MHz에서 슬랙 $+0.18$나노초였다(\ref{sec:timing}절).
배치와 배선을 해 보니 결과가 달랐다. 배치 뒤의 슬랙이 $-3.4$나노초, 배선 중에는 $-5.2$나노초였다.
가장 느린 경로는 예상대로 전파 회로였고 약 15나노초가 걸렸는데, 그 대부분이 LUT가 아니라 \emph{배선}의
지연이었다. 676비트의 상태가 메뉴의 선 16개의 양쪽 끝, 곧 32곳의 선택 회로로 퍼져 나가야 하므로 칩 위에서
먼 거리를 가는 선이 많다. 합성 단계의 추정이 이 거리를 낮게 어림한 것이다.

고칠 방법은 \ref{sec:timing}절의 셋이다. 전파 회로는 고리이므로 파이프라인으로 나누면 걸음마다 두 클럭이
들어 얻는 것이 없다. 구조를 더 바꾸는 것은 가능하지만 큰 일이다. 이 프로젝트의 최종 목표는 기계식 실물이고
FPGA는 세 번째 전시물이자 대조용 기준 장치이므로, 셋째 방법을 골랐다. 시스템 클럭을 50MHz로 내렸다.
보드의 100MHz 클럭을 칩 안의 클럭 생성기로 절반으로 나눈다.
\end{lesson}

50MHz에서의 결과가 표~\ref{tab:g-e-impl}이다. 대상은 Nexys Video 보드의 XC7A200T(LUT 약 13만 4천 개)이고
도구는 Vivado다.

\begin{table}[ht]
\centering
\caption{배치와 배선을 마친 결과(50MHz, 2026년 11월 15일). 슬랙은 가장 나쁜 경로의 값이다.}
\label{tab:g-e-impl}
\small
\begin{tabular}{rrrr}
\toprule
엔진 수 & LUT & 슬랙 & 빌드 시간 \\
\midrule
1 & 43,100 (32.2\%) & $+2.40$ ns & 52분 \\
2 & 83,257 (62.2\%) & $+0.13$ ns & 3시간 52분 \\
\bottomrule
\end{tabular}
\end{table}

엔진이 한 대일 때에는 여유가 2.4나노초인데 두 대가 되면 0.13나노초로 줄고 빌드 시간이 네 배 넘게
늘었다. 회로는 똑같은 것 두 벌인데도 그렇다. \ref{ch:pipeline}장에서 말한 대로 칩이 찰수록 부품을 가까이
놓을 자리가 없어 배선이 길어지기 때문이다. 세 대이면 LUT의 약 93\%를 쓰게 되어 배선이 현실적이지 않다.
그래서 이 보드의 엔진은 두 대로 정했다. ``들어간다''(자원)와 ``제 속도로 돈다''(타이밍)가 서로 다른 한계라는
것, 그리고 자원의 한계보다 타이밍의 한계가 먼저 온다는 것을 보여 주는 예다.

속도는 위치당 약 15클럭으로 잡으면 엔진 한 대가 초당 약 330만 위치다. 로터 순서 하나(17,576 위치)가
약 5밀리초다. 실제 봄베는 같은 일에 약 10분이 걸렸다. 한편 Rust 코어는 PC의 스레드 24개로 초당 약
900만에서 1,000만 위치를 낸다. 엔진 두 대를 합쳐도 PC보다 빠르지 않다. 이 엔진의 목적은 속도 기록이 아니라
같은 계산을 움직이는 부품 없이 회로로 해 보이는 것이다. 이 속도는 시뮬레이션의 클럭 수에서 계산한 값이고
보드에서 잰 값이 아니다.

\begin{todobox}{보드에서의 실측과 전시용 표시}
보드를 확보한 뒤의 실측 속도와 호스트를 통한 Rust 결과와의 대조, PYNQ 보드의 빌드 결과. 전시물로서 내부
상태(676가닥의 선, 드럼 위치, 정지)를 눈에 보이게 하는 표시 장치와 느린 시연 모드의 구조.
\end{todobox}

\section{G부의 정리}

\begin{center}
\small
\begin{tabularx}{\linewidth}{>{\raggedright\arraybackslash}p{38mm}>{\raggedright\arraybackslash}X>{\raggedright\arraybackslash}X}
\toprule
실제 봄베 & 기준 RTL & FPGA 엔진 \\
\midrule
드럼을 골라 끼운다 & 로터 번호로 상수 표를 매번 고른다 & 구동 시작 때 표를 레지스터로 옮긴다 (\code{LOAD}) \\
느린 드럼은 26위치 동안 서 있다 & 이용하지 않는다 & 가운데 로터가 바뀔 때만 표를 다시 만든다 (\code{INNER}, \code{SROWS}) \\
드럼 한 벌 & 스크램블러 조합 논리, 선마다 한 벌 & 행 메모리의 사본에서 읽은 \code{perm\_q} \\
유닛마다 드럼을 크립 위치만큼 어긋나게 끼운다 & 위치에 크립 위치를 더한다 & 읽는 주소에 크립 위치를 더한다 \\
676가닥의 전선 & 676비트 한 줄 \code{live\_q} & 26비트 $\times$ 26 \code{cab\_q} \\
메뉴의 케이블 배선 & 정수 번호로 비트를 고른다 & 케이블 선택, 치환, 원핫으로 모으기 \\
감지 릴레이 26개 & \code{test\_q != FULL} & 같다. 고정점은 레지스터 둘을 비교 \\
정지하고 조작자를 기다린다 & \code{EMIT} 상태와 valid/ready & FIFO에 넣고 계속 간다 \\
막사의 봄베 여러 대 & --- & \code{bombe\_array}의 엔진 $N$대 \\
\bottomrule
\end{tabularx}
\end{center}

\begin{keybox}[title={G부에서 쓴 방법들}]
\begin{itemize}
\item 입력이 드물게 바뀌는 계산은 따로 떼어, 작은 회로로 천천히 계산하고 결과를 저장한다.
\item 실행 중에 고르지만 고른 뒤에는 바뀌지 않는 것은 시작할 때 한 번 풀어 둔다(배선표, 원핫).
\item 파이프라인에는 값과 함께 꼬리표와 유효 비트를 흘린다.
\item 한 줄로 차례로 할지 여러 줄로 한꺼번에 할지는 계산의 양과 횟수로 정한다.
\item 메모리는 좁지만 읽는 단자가 적다. 동시에 여러 곳을 읽어야 하면 사본을 둔다.
\item 가장 긴 경로 위에서는 일을 덜어 낸다. 판정은 레지스터에 담긴 값으로 한다.
\item 구조가 달라질수록 기준 모델과의 대조에 기댄다.
\end{itemize}
\end{keybox}
